\naslov{Pričakovana vrednost in varianca slučajnih vektorjev}

\begin{definicija}
  Za slučajni vektor $\sv{X} = (X_1, \ldots, X_n)$ definiramo $E(\sv{X}) =
  (E(X_1), \ldots, E(X_n))$.
\end{definicija}

\begin{trditev}
  Za deterministično matriko $A$ primerne velikosti in slučajni vektor $\sv{X}$
  velja $E(A \sv{X}) = A E(\sv{X})$.
\end{trditev}

\begin{posledica}
  Če je $\sv{u}$ determinističen, $\sv{X}$ pa slučajen vektor, je $E(\sv{u}
  \cdot \sv{X}) = \sv{u} \cdot E(\sv{X})$.
\end{posledica}

\begin{definicija}
  Za slučajno matriko $M$ definiramo $E(M) = [E(M_{ij})]_{ij}$.
\end{definicija}

\begin{trditev}
  Naj bo $M$ slučajna matrika.
  Potem je $E(M^T) = E(M)^T$.
  Za poljubni deterministični matriki $A$ in $B$ primerne velikosti velja $E(A M
  B) = A E(M) B$.
\end{trditev}

\begin{trditev}
  Za poljubna primerna slučajna vektorja $\sv{X}$ in $\sv{Y}$ ter za poljubni
  primerni slučajni matriki $M$ in $N$ velja $E(\sv{X} + \sv{Y}) = E(\sv{X}) +
  E(\sv{Y})$ in $E(M + N) = E(M) + E(N)$.
\end{trditev}

\begin{definicija}
  \pojem{Kovariančna matrika} med slučajnimi vektorji $\sv{X}$ in $\sv{Y}$ je
  \[
	\cov(\sv{X}, \sv{Y}) =
	\begin{bmatrix}
	  \cov(X_1, Y_1) & \cdots & \cov(X_1, Y_n) \\
	  \vdots & \ddots & \vdots \\
	  \cov(X_n, Y_1) & \cdots & \cov(X_n, Y_n)
	\end{bmatrix}.
  \]
\end{definicija}

\begin{trditev}
  $\cov(\sv{X}, \sv{Y}) = E((\sv{X} - E(\sv{X}))(\sv{Y} - E(\sv{Y}))^T) =
  E(\sv{X} \sv{Y}^T) - E(\sv{X}) E(\sv{Y})^T$.
\end{trditev}

\vprasanje{Kako izraziš kovariančno matriko?}

Izračunamo lahko
\[
  \cov(A \sv{X}, B \sv{Y})
  = A \cov(\sv{X}, \sv{Y}) B^T,
\]
iz česar za vektor $\sv{u}$ dobimo
\[
  0 \le \var(\sv{u} \cdot \sv{X})
  = \sv{u}^T \var(\sv{X}) \sv{u}
  = \sk{\var(\sv{X}) \sv{u}, \sv{u}},
\]
torej je kovariančna matrika slučajnega vektorja pozitivno semidefinitna.

\begin{trditev}
  Pozitivno semidefinitna matrika je kovariančna matrika slučajnega vektorja.
\end{trditev}

\begin{proof}
  Naj bo $\Sigma$ pozitivno semidefinitna matrika.
  Vzemimo $\sv{Z} = (Z_1, \ldots, Z_n)$, kjer so $Z_i$ neodvisne in porazdeljene
  standardno normalno.
  Velja $\var(\sv{Z}) = I_n$.
  Za $\Sigma$ obstaja razcep Choleskega $\Sigma = V V^T$.
  Potem je $\cov(A \sv{Z}, A \sv{Z}) = A A^T = \Sigma$.
\end{proof}

\vprasanje{Pokaži, da je matrika enaka kovariančni matriki slučajnega vektorja
  natanko tedaj, ko je pozitivno semidefinitna.}

\begin{definicija}
  Slučajna vektorja $\sv{X}$ in $\sv{Y}$ sta \pojem{nekorelirana}, če je
  $\cov(\sv{X}, \sv{Y}) = 0$.
\end{definicija}
